\documentclass[11pt,a4paper]{article}

\usepackage[T1]{fontenc}
\usepackage[utf8]{inputenc}
\usepackage[english]{babel}
\usepackage{lmodern}
\usepackage{microtype}
\usepackage[a4paper,margin=20mm,headheight=15pt]{geometry}
\usepackage{graphicx}
\usepackage{booktabs}
\usepackage{tabularx}
\usepackage{longtable}
\usepackage{array}
\usepackage{xcolor}
\usepackage{amsmath}
\usepackage{caption}
\usepackage{float}
\usepackage{enumitem}
\usepackage{fancyhdr}
\usepackage[hidelinks]{hyperref}
\usepackage{xurl}

\graphicspath{{figures/}}
\definecolor{controlblue}{HTML}{2563EB}
\definecolor{acceptedgreen}{HTML}{16A34A}
\definecolor{rejectedred}{HTML}{DC2626}
\definecolor{amber}{HTML}{D97706}
\definecolor{slate}{HTML}{334155}
\definecolor{lightgray}{HTML}{F1F5F9}

\newcommand{\dstwr}{DS-TWR}
\newcommand{\code}[1]{\texttt{\detokenize{#1}}}
\newcolumntype{Y}{>{\raggedright\arraybackslash}X}

\pagestyle{fancy}
\fancyhf{}
\lhead{\small Native DS-TWR speed-limit revalidation}
\rhead{\small July 31, 2026}
\cfoot{\thepage}

\hypersetup{
  pdftitle={Native DS-TWR speed-limit revalidation},
  pdfauthor={UWB ESP-IDF technical report},
  pdfsubject={Fresh static precision, raw-range integrity, and speed-boundary study},
  pdfkeywords={UWB, DS-TWR, DW3000, ESP32-S3, ranging, precision, integrity}
}

\begin{document}
\pagenumbering{gobble}

\begin{titlepage}
  \centering
  \vspace*{15mm}
  {\Huge\bfseries Native \dstwr{} speed-limit revalidation\par}
  \vspace{4mm}
  {\Large Fresh throughput, precision, and raw-range integrity study\par}
  \vspace{15mm}

  \begin{tabular}{rl}
    Test date: & July 31, 2026, EEST\\
    Platform: & 5 ESP32-S3 + DW3000 modules, Raspberry Pi 5\\
    Branch: & \texttt{experiment/passive-ds-twr-pipeline}\\
    Dashboard source: & \texttt{8244a3432f68}\\
    Active firmware identity: & \texttt{54a586c-dirty}\\
    Configuration: & M1 tag; M2, M3, M4, M5 anchors\\
    Radio: & UWB channel 5\\
    Environment: & unchanged static indoor setup\\
  \end{tabular}

  \vfill
  \colorbox{lightgray}{%
    \begin{minipage}{0.88\textwidth}
      \vspace{3mm}
      \textbf{Key result.}
      The \textbf{60 ms} profile is the fastest setting that completed the
      decisive 180 s validation with \textbf{zero impossible raw ranges}.
      It delivered \textbf{11.63 coherent independent frames/s} and
      \textbf{61.96 rolling display updates/s}, with
      \textbf{1.38 cm precision RMS} and \textbf{2.39 cm CEP95}.
      Relative to the fresh 410 ms control, this is a
      \textbf{6.75$\times$ coherent-frame speedup} for only
      \textbf{+1.6\% RMS} and \textbf{+3.9\% CEP95}.
      The otherwise excellent-looking 59 ms profile was rejected because it
      produced a physically impossible $-5.301$ m raw range.
      \vspace{3mm}
    \end{minipage}}
  \vfill
  {\large Reproducible updated field report\par}
  \vspace{5mm}
\end{titlepage}

\pagenumbering{roman}
\tableofcontents
\clearpage
\pagenumbering{arabic}

\section{Executive summary}

This report revalidates the native three-message \dstwr{} speed limit after
restoring the low-noise ranging behavior and capturing a new slow control on
the same physical setup. The exchange remains strictly
\texttt{POLL--RESP--FINAL}. No temporal filter, moving average, position
smoothing, or additional radio message is used in the offline comparison.

The study contains 17 capture blocks, 1,860.8 s of data, and 156,935 raw
tag-to-anchor ranges. It first repeats the dashboard profiles from 410 ms down
to 13 ms, then narrows the transition using 49, 57, 58, 59, 60, 61, and
64 ms candidates. Decisive profiles use 180 s blocks; short 60 s screens are
used only to reject obviously unsafe settings or select the next candidate.

\begin{table}[H]
  \centering
  \caption{Main fresh results. Precision is measured around each unfiltered
  coherent-position series' own mean.}
  \small
  \begin{tabular}{lrrrrrr}
    \toprule
    Profile & Coherent & Rolling & Delivery & RMS & CEP95 & Raw anomalies\\
    \midrule
    410 ms control & 1.72 Hz & 9.16 Hz & 94.11\% & 1.36 cm & 2.30 cm & 0\\
    64 ms reserve & 10.92 Hz & 58.14 Hz & 94.03\% & 1.38 cm & 2.41 cm & 0\\
    \textcolor{acceptedgreen}{\textbf{60 ms recommended}}
      & \textbf{11.63 Hz} & \textbf{61.96 Hz} & \textbf{93.97\%}
      & \textbf{1.38 cm} & \textbf{2.39 cm} & \textbf{0}\\
    \textcolor{rejectedred}{59 ms rejected}
      & 11.82 Hz & 62.78 Hz & 93.75\% & 1.31 cm & 2.20 cm & 1\\
    \textcolor{rejectedred}{58 ms rejected}
      & 12.00 Hz & 63.48 Hz & 93.32\% & 1.35 cm & 2.30 cm & 2\\
    \bottomrule
  \end{tabular}
\end{table}

\paragraph{Decision.}
Use 60 ms as the fastest currently validated precision profile. Use 64 ms
when an additional millisecond of reserve per anchor slot is preferred.
Retain 410 ms only as a slow control. Treat every profile below 60 ms as an
experimental stress test: their central position distributions can remain
tight while rare timestamp failures generate impossible raw distances.

\section{Why the recommendation changed}

The July 26 report recommended 17 ms and treated 13 ms as the measured
maximum. That report faithfully described the complete-position and timing
statistics available in its original capture. The present result does not
erase that historical data; it changes the operating recommendation for the
current firmware and validation criterion.

Three differences matter:

\begin{enumerate}[leftmargin=8mm]
  \item the low-noise native \dstwr{} implementation was recovered and a fresh
  410 ms control was captured on the unchanged current setup;
  \item the current dashboard maintains live anchor geometry and distinguishes
  rolling estimates from strict coherent four-anchor frames;
  \item most importantly, every raw range is now audited before solver or
  frame-assembly rejection. Negative ranges and deviations greater than
  10 cm from the per-anchor median are counted even when they do not visibly
  disturb the position trail.
\end{enumerate}

The new audit found impossible raw measurements at 33, 29, 25, 21, 17, and
13 ms. Their coherent RMS and CEP95 still look comparable to the control.
Position precision alone is therefore insufficient to qualify the radio
timing. The updated recommendation requires both tight central precision and
clean raw timestamp integrity.

\section{System and protocol}

\subsection{Native three-frame DS-TWR}

For each anchor, the tag sends \texttt{POLL}, the anchor returns
\texttt{RESP}, and the tag sends \texttt{FINAL}. The anchor combines its
hardware receive/transmit timestamps with the tag timestamps embedded in the
final message and estimates propagation time as

\begin{equation}
  \widehat{T}_{\mathrm{prop}} =
  \frac{R_aR_b-D_aD_b}{R_a+D_a+R_b+D_b},
\end{equation}

where $R_a,R_b$ are round-trip intervals and $D_a,D_b$ are reply delays.
Antenna-delay calibration and the instantaneous carrier-integrator correction
remain enabled. They correct physical timestamp behavior and are not temporal
position filters.

\subsection{Frame and geometry scheduler}

The tag visits anchors A2, A3, A4, and A5 sequentially. The base tag-ranging
frame is

\begin{equation}
  T_{\mathrm{base}} = 4T_{\mathrm{slot}} + T_{\mathrm{gap}}.
\end{equation}

The current implementation also inserts one rotating anchor-to-anchor
geometry-maintenance slot every four base frames. With a 1 ms command delay,
the average scheduler period is approximately

\begin{equation}
  T_{\mathrm{avg}} =
  T_{\mathrm{base}} +
  \frac{T_{\mathrm{slot}}+1\ \mathrm{ms}}{4}.
\end{equation}

This maintenance exchange keeps the displayed physical geometry live and is
included in measured throughput. It is one reason the measured strict
coherent-frame rate is lower than $1/T_{\mathrm{base}}$.

\subsection{Coherent frames and rolling updates}

A coherent independent position uses four ranges assigned to one complete
anchor round. A rolling position is recomputed whenever a new range arrives,
using the latest still-fresh ranges from the other anchors. Rolling updates
provide low display latency, but adjacent estimates share measurements and
must not be interpreted as statistically independent frames.

At 60 ms, the dashboard can produce 61.96 rolling updates/s while the strict
offline assembler identifies 11.63 complete independent frames/s. Both rates
are reported so display responsiveness is not confused with independent
measurement throughput.

\section{Experimental method}

\subsection{Geometry and unchanged setup}

All captures use the same cached measured geometry:

\begin{center}
\begin{tabular}{lrr}
  \toprule
  Anchor & $x$ [m] & $y$ [m]\\
  \midrule
  A2 & 0.000 & 0.000\\
  A3 & 0.000 & 3.067\\
  A4 & 2.971 & $-0.071$\\
  A5 & 3.196 & 2.974\\
  \bottomrule
\end{tabular}
\end{center}

The stored tag reference is the anchor centroid,
$(1.54175,1.49250)$ m. Because this coordinate is derived from cached
geometry rather than an independent survey of the tag location, the report
uses dispersion around each capture mean as its primary precision metric.
Absolute bias fields are retained in \path{summary.json} but are not presented
as independently verified absolute accuracy.

Active antenna delays remained M1=16363, M2=16368, M3=16362, M4=16371, and
M5=16371. No module was physically moved during the sweep. Runtime profile
changes were applied to all five modules in parallel, followed by reboot and
warm-up.

\subsection{Fresh control and boundary search}

The sequence was:

\begin{enumerate}[leftmargin=8mm]
  \item capture a new 180 s 410 ms precision control;
  \item validate the existing 64, 33, and 29 ms profiles for 180 s and screen
  25, 21, 17, and the true 13 ms profile for 60 s;
  \item test intermediate 49, 57, 58, 60, and 61 ms candidates for 60 s;
  \item perform decisive 180 s validations at 58, 59, and 60 ms.
\end{enumerate}

The short 58 ms screen was clean, but its 180 s validation produced two
impossible ranges. Likewise, the 29 ms screen was clean and its long
validation failed. These two cases demonstrate why short visual screening
cannot establish reliability.

\subsection{Acceptance criteria}

A candidate is accepted only if:

\begin{itemize}[leftmargin=7mm]
  \item the decisive 180 s block contains no negative raw distance;
  \item no raw distance differs by more than 10 cm from that anchor's capture
  median;
  \item coherent RMS and CEP95 remain within 20\% of the 410 ms control;
  \item all five modules remain online with stable sequence delivery.
\end{itemize}

The 10 cm threshold is an integrity alarm, not the claimed localization
accuracy. Every rejected boundary event was several metres beyond that limit.
No range is made acceptable merely because a later solver ignores it.

\subsection{Metrics}

The offline reconstruction uses unweighted all-anchor linear trilateration and
no temporal filtering. Precision RMS is

\begin{equation}
  \mathrm{RMS}_{\mathrm{precision}} =
  \sqrt{\frac{1}{N}\sum_{i=1}^{N}
  \left\lVert \mathbf{p}_i-\overline{\mathbf{p}}\right\rVert^2}.
\end{equation}

CEP50 and CEP95 are the corresponding 50th and 95th percentiles of radial
deviation from the capture mean. Sequence delivery is the fraction of observed
16-bit range sequences within the capture span. Raw integrity is calculated
directly from every \texttt{UWB\_RANGING result} record.

\section{Results}

\subsection{Update rate and scheduler yield}

\begin{figure}[H]
  \centering
  \includegraphics[width=\textwidth]{01_rate_delivery.pdf}
  \caption{Measured coherent and rolling update rates, plus the yield of each
  stream relative to its ideal average scheduler rate. An asterisk marks a
  profile that failed raw-integrity validation. Colors retain the metric
  meanings shown in the legend.}
\end{figure}

The 60 ms profile yields a 6.75$\times$ coherent-frame speedup and a
6.77$\times$ rolling-update speedup over the 410 ms control. Sequence delivery
changes by only $-0.14$ percentage points. The lower timing profiles raise
apparent update rate, but raw integrity fails. In the lower panel, coherent
yield is the observed independent-frame rate divided by the ideal
geometry-adjusted frame rate. Rolling yield uses four ideal display updates
per complete tag-ranging frame. The difference quantifies how shared
fresh-range updates remain responsive even when strict four-range frames are
incomplete.

\subsection{Central precision}

\begin{figure}[H]
  \centering
  \includegraphics[width=0.94\textwidth]{02_central_precision.pdf}
  \caption{Unfiltered coherent-position dispersion. The rejected profiles
  are marked with an asterisk but still have deceptively good RMS and CEP95.
  Colors retain the RMS/CEP95 meanings shown in the legend.}
\end{figure}

The 60 ms coherent RMS is 1.3766 cm, compared with 1.3555 cm at 410 ms:
+0.0211 cm or +1.56\%. CEP95 changes from 2.2987 to 2.3892 cm:
+0.0905 cm or +3.94\%. These are small differences relative to ordinary
between-capture variation.

\begin{figure}[H]
  \centering
  \includegraphics[width=0.91\textwidth]{05_precision_cdf.pdf}
  \caption{Empirical CDF of radial coherent-position deviation from each
  capture mean. The accepted 410, 64, and 60 ms distributions nearly overlap.}
\end{figure}

The CDF confirms that the recommended 60 ms setting preserves the complete
shape of the static precision distribution rather than only its RMS summary.

\subsection{Raw-range integrity}

\begin{figure}[H]
  \centering
  \includegraphics[width=0.95\textwidth]{03_raw_anomalies.pdf}
  \caption{Impossible measurements per 10,000 raw ranges, counted before any
  frame or solver rejection.}
\end{figure}

The accepted 410, 64, and 60 ms captures contain no raw anomaly. Every
existing dashboard profile below 60 ms produces at least one impossible
measurement in its fresh capture. At 17 ms, 13 negative ranges and 15 total
spikes appear in 11,467 observations; at 13 ms, the counts rise to 20 and 22
in 14,303 observations.

\begin{figure}[H]
  \centering
  \includegraphics[width=\textwidth]{04_boundary_integrity.pdf}
  \caption{All A2 raw deviations during the decisive 58, 59, and 60 ms
  validations. The symlog axis preserves centimetre-scale noise and
  metre-scale failures in the same figure.}
\end{figure}

The boundary result is unusually sharp. The 58 ms profile produced two
negative A2 ranges, the 59 ms profile produced one, and the otherwise
identical 60 ms profile produced none. All three use 14 ms anchor slots,
8 ms RX timeout, and 2+2 ms response/final delays. Only the frame gap changes:
2, 3, and 4 ms respectively.

\subsection{Accepted per-anchor stability}

\begin{figure}[H]
  \centering
  \includegraphics[width=0.90\textwidth]{06_anchor_stability.pdf}
  \caption{Raw range standard deviation for the three accepted profiles.}
\end{figure}

The accepted profiles retain approximately 1--2 cm raw standard deviation per
anchor. The 60 ms values are 1.19 cm (A2), 1.12 cm (A3), 1.23 cm (A4), and
1.79 cm (A5). There is no broad noise inflation as the frame period falls
from 410 to 60 ms.

\section{Interpretation and recommendation}

\subsection{Observed limiting mechanism}

The data do not show a gradual accuracy-versus-speed tradeoff near the
boundary. Central RMS is actually slightly lower at 59 ms than at 60 ms.
Instead, rare metre-scale failures appear while the ordinary range population
remains centimetre-stable.

Because 58, 59, and 60 ms use the same 14 ms slot and exchange delays, while
only the inter-frame gap changes, the evidence strongly implicates insufficient
recovery margin between rounds. Every decisive boundary failure is emitted by
A2, the first anchor in the configured list. This pattern is consistent with a
round-boundary timestamp or radio-state race, but the present field data do
not prove the exact firmware instruction or DW3000 state responsible.

The correct engineering response is therefore to retain a clean timing margin,
not to add a position filter that hides the symptom. Instrumenting the
round transition and validating an explicit timestamp sanity check remain
useful future firmware tasks.

\subsection{Recommended profiles}

\paragraph{Default: 60 ms precision-speed limit.}
Use 14 ms slots, 4 ms frame gap, 8 ms RX timeout, 2 ms response delay, 2 ms
final delay, 100 ms anchor RX slice, and 500 UUS auto-RX delay. This is the
fastest profile with a clean 180 s raw audit.

\paragraph{Reserve: 64 ms.}
Use 15 ms slots and the same 4/8/2/2 ms gap/timeout/response/final values.
It loses about 6\% coherent throughput relative to 60 ms but also passed its
180 s capture with zero raw anomalies.

\paragraph{Control: 410 ms.}
Use 100 ms slots, 10 ms gap, 90 ms timeout, and 20+20 ms delayed
transmissions only when reproducing the slow precision baseline. It is no
longer justified as the normal operating profile.

\paragraph{Experimental profiles below 60 ms.}
They remain available in the dashboard for controlled investigation, but are
now marked rejected or experimental. Their high rolling-update numbers must
not be presented as validated reliable positioning throughput.

\section{Limits of inference and next tests}

This is a strict static revalidation at one favorable point, not a complete
localization qualification. The 60 ms decision applies to the tested modules,
firmware state, channel, anchor order, geometry, and host load. It does not yet
establish:

\begin{itemize}[leftmargin=7mm]
  \item multi-hour probability of a rare boundary failure;
  \item performance while the tag translates or rotates;
  \item behavior near or outside the anchor perimeter;
  \item NLOS, human blockage, or severe multipath performance;
  \item sensitivity to Wi-Fi/dashboard load, temperature, or supply voltage;
  \item behavior with a different first anchor or additional modules.
\end{itemize}

The next validation should run 60 and 64 ms for at least one hour, alternate
their order, and repeat with A2 moved away from the first slot. A controlled
dynamic trajectory should then compare coherent latency and precision. If
sub-60 ms operation remains important, add per-stage round-transition
timestamps and an explicit hardware-timestamp validity guard before changing
the estimator or adding filters.

\clearpage
\appendix

\section{Complete profile table}

\begin{longtable}{rrrrrrrrr}
  \caption{Fresh aggregate results. Durations are 180 s unless marked
  60 s screen.}\\
  \toprule
  Frame & Duration & Coherent & Rolling & Delivery & RMS & CEP95 & Neg. & Spikes\\
  \midrule
  \endfirsthead
  \toprule
  Frame & Duration & Coherent & Rolling & Delivery & RMS & CEP95 & Neg. & Spikes\\
  \midrule
  \endhead
  410 & 180 s & 1.72 & 9.16 & 94.11\% & 1.36 & 2.30 & 0 & 0\\
  64 & 180 s & 10.92 & 58.14 & 94.03\% & 1.38 & 2.41 & 0 & 0\\
  60 & 180 s & 11.63 & 61.96 & 93.97\% & 1.38 & 2.39 & 0 & 0\\
  59 & 180 s & 11.82 & 62.78 & 93.75\% & 1.31 & 2.20 & 1 & 1\\
  58 & 180 s & 12.00 & 63.48 & 93.32\% & 1.35 & 2.30 & 2 & 2\\
  57 & 60 s & 12.21 & 64.37 & 93.02\% & 1.38 & 2.42 & 1 & 1\\
  49 & 60 s & 17.61 & 74.26 & 92.46\% & 1.30 & 2.25 & 1 & 1\\
  33 & 180 s & 20.42 & 107.21 & 90.61\% & 1.33 & 2.33 & 7 & 9\\
  29 & 180 s & 23.37 & 119.09 & 88.92\% & 1.43 & 2.46 & 13 & 19\\
  25 & 60 s & 27.09 & 137.28 & 89.28\% & 1.43 & 2.47 & 2 & 2\\
  21 & 60 s & 33.34 & 157.85 & 87.90\% & 1.43 & 2.53 & 4 & 4\\
  17 & 60 s & 39.77 & 191.04 & 87.14\% & 1.34 & 2.35 & 13 & 15\\
  13 & 60 s & 44.99 & 238.12 & 85.30\% & 1.38 & 2.36 & 20 & 22\\
  \bottomrule
\end{longtable}

Coherent and rolling columns are updates/s. RMS and CEP95 are centimetres.
Spikes count all raw distances more than 10 cm from the corresponding
per-anchor median and include the negative observations.

\section{Timing profiles}

\begin{table}[H]
  \centering
  \small
  \begin{tabular}{rrrrrr}
    \toprule
    Frame & Slot & Gap & RX timeout & RESP delay & FINAL delay\\
    \midrule
    410 & 100 & 10 & 90 & 20 & 20\\
    64 & 15 & 4 & 8 & 2 & 2\\
    60 & 14 & 4 & 8 & 2 & 2\\
    59 & 14 & 3 & 8 & 2 & 2\\
    58 & 14 & 2 & 8 & 2 & 2\\
    57 & 14 & 1 & 8 & 2 & 2\\
    49 & 12 & 1 & 8 & 2 & 2\\
    33 & 8 & 1 & 5 & 2 & 2\\
    29 & 7 & 1 & 5 & 2 & 2\\
    25 & 6 & 1 & 5 & 2 & 2\\
    21 & 5 & 1 & 5 & 2 & 2\\
    17 & 4 & 1 & 4 & 1 & 1\\
    13 & 3 & 1 & 3 & 1 & 1\\
    \bottomrule
  \end{tabular}
  \caption{Runtime values in milliseconds. RX slice is 100 ms and automatic
  RX delay is 500 UUS for every row.}
\end{table}

\section{Files and reproduction}

\small
\begin{tabularx}{\textwidth}{p{62mm}Y}
  \toprule
  File/directory & Contents\\
  \midrule
  \path{README.md} & fresh 410 ms control identity and precision reference\\
  \path{PROFILE_REVALIDATION.md} & concise findings and boundary decision\\
  \path{profile_comparison.csv} & machine-readable aggregate table\\
  \path{candidate_*ms_validation/data/*.jsonl} & retained RAW event streams
  for the decisive 58/59/60 ms boundary\\
  \path{*/data/*.metadata.json} & full initial/final module status and runtime
  configuration\\
  \path{*/positions.csv} & unfiltered coherent and rolling reconstructions\\
  \path{*/summary.json} & complete per-capture metrics\\
  \path{*/raw_quality.json} & per-anchor raw integrity audit\\
  \path{*/SHA256SUMS} & historical hashes for the original capture files\\
  \path{figures/} & report plots in vector PDF and PNG\\
  \bottomrule
\end{tabularx}

\paragraph{Reproduction}
The repository uses an aggressive retention policy: the current figure
generator can rebuild the published plots from the three retained boundary
event streams and the derived per-profile artifacts. Re-auditing every
profile from its original RAW stream is no longer possible.
\begin{verbatim}
cd /home/pi/Documents/UWB
python3 tools/uwb_native_ds_speed_revalidation_report.py
cd reports/native_ds_speed_revalidation_20260731
latexmk -pdf -interaction=nonstopmode -halt-on-error report.tex
\end{verbatim}

\begin{thebibliography}{9}

\bibitem{qorvo-api}
Qorvo, \emph{DW3xxx Device Driver API Guide}, double-sided two-way-ranging
examples, 2025.
\url{https://forum.qorvo.com/uploads/short-url/xD3TlXKvkujjdUaJWXv2b4E7GQN.pdf}.

\bibitem{aps013}
Decawave, \emph{APS013: DW1000 and Two-Way Ranging}, rev. 2.3,
Poll--Response--Final and the asymmetric DS-TWR formula.
\url{https://forum.qorvo.com/uploads/short-url/x34DrF7EW5fQP9wY3aNESqPKz8z.pdf}.

\bibitem{firmware}
Local repository \emph{UWB Module ESP-IDF}, branch
\texttt{experiment/passive-ds-twr-pipeline}, dashboard source
\texttt{8244a3432f68caf23d604e475cbe718698c4ab2f}, active module firmware
identity \texttt{54a586c-dirty}, field state from July 31, 2026.

\end{thebibliography}

\end{document}
